\begin{theorem}[Markov’s inequality]
Let $(\Omega,\mathcal F,\mathbb P)$ be a probability space and let $X:\Omega\to[0,\infty]$ be a non‑negative random variable.  
For any $a>0$,
\[
\mathbb P\bigl(X\ge a\bigr)\;\le\;\frac{\mathbb E[X]}{a}.
\]
Equality holds iff $X$ takes only the two values $0$ and $a$ almost surely (equivalently $X=a\,\mathbf 1_{\{X\ge a\}}$ a.s.).
\end{theorem}

\begin{proof}
\textbf{1.  Assumptions and degenerate case.}  
If $\mathbb P(X\ge a)=0$, the left–hand side of the desired inequality is $0$, while $\mathbb E[X]\ge0$ and $a>0$ imply $\mathbb E[X]/a\ge0$. Hence the inequality holds trivially, and we may assume $\mathbb P(X\ge a)>0$.

\medskip\noindent\textbf{2.  Indicator of the event $\{X\ge a\}$.}  
Define
\[
\mathbf 1_{\{X\ge a\}}(\omega)=
\begin{cases}
1, & X(\omega)\ge a,\\[4pt]
0, & X(\omega)< a .
\end{cases}
\]

\medskip\noindent\textbf{3.  Pointwise bound.}  
For every $\omega\in\Omega$ we have
\[
a\,\mathbf 1_{\{X\ge a\}}(\omega)\le X(\omega).
\tag{1}
\]
Indeed, if $X(\omega)\ge a$ then the left‑hand side equals $a\le X(\omega)$; if $X(\omega)<a$ the left‑hand side is $0\le X(\omega)$ by non‑negativity of $X$.

\medskip\noindent\textbf{4.  Taking expectations.}  
The expectation operator is monotone, so applying it to \eqref{1} gives
\[
a\,\mathbb E\bigl[\mathbf 1_{\{X\ge a\}}\bigr]\le\mathbb E[X].
\tag{2}
\]

\medskip\noindent\textbf{5.  Expectation of an indicator.}  
By definition,
\[
\mathbb E\bigl[\mathbf 1_{\{X\ge a\}}\bigr]=\mathbb P(X\ge a).
\tag{3}
\]
Substituting \eqref{3} into \eqref{2} and dividing by the positive constant $a$ yields
\[
\mathbb P(X\ge a)\le\frac{\mathbb E[X]}{a},
\]
which is Markov’s inequality.

\medskip\noindent\textbf{6.  Equality condition.}  
Equality in the final inequality can occur only if every inequality used above is an equality almost surely.

\begin{itemize}
\item Equality in the monotonicity step \eqref{2} forces the pointwise bound \eqref{1} to hold with equality $\mathbb P$‑a.s., i.e.
\[
a\,\mathbf 1_{\{X\ge a\}}(\omega)=X(\omega)\quad\text{for $\mathbb P$‑a.e. }\omega .
\]
\item This identity means $X(\omega)=a$ on the set $\{X\ge a\}$ and $X(\omega)=0$ on its complement.
\end{itemize}
Thus $X$ must be of the form $X=a\,\mathbf 1_{\{X\ge a\}}$ a.s., i.e. $X$ takes only the values $0$ and $a$. Conversely, for such a two‑point distribution we have $\mathbb E[X]=a\,\mathbb P(X=a)$, so $\mathbb P(X\ge a)=\mathbb E[X]/a$ and equality holds.

\medskip\noindent\textbf{7.  Conclusion.}  
For every non‑negative random variable $X$ and every $a>0$,
\[
\mathbb P\bigl(X\ge a\bigr)\le\frac{\mathbb E[X]}{a},
\]
with equality precisely when $X$ assumes only the values $0$ and $a$ almost surely. \qed
\end{proof}